\chapter{Technical notes}\label{ch:technical}

This chapter describes how \nwrem{} computes its estimates and how fast it is.

\section{The structure of a fit}\label{sec:structure}

A fit has four steps. (1) The events are read from the network, checked and sorted by time (events with the same time keep the order of the rows). (2) The model, that is the events, the covariates, the risk set and the windows, is written once into a data frame and handed to the plugin; the frame is dropped at the end, also after an error. (3) The plugin checks for effects that are exact functions of others, by one evaluation at $\theta=0$ (Section~\ref{sec:collinear}); the evaluation at $\theta=0$ serves at the same time as the log likelihood of the null model and as the first Newton step. (4) A damped Newton iteration maximizes the log likelihood, calling the plugin once per evaluation with the current coefficients; each call returns the log likelihood, the gradient and the exact Hessian.

\paragraph{Newton iteration.} The log likelihood is concave, so the Newton direction $d=(-H)^{-1}g$ is an ascent direction. A full step is taken and halved (up to 30 times) until the log likelihood does not fall. The iteration stops when the Newton decrement $g^{\top}(-H)^{-1}g$ is below $\max(10^{-12},10^{-13}|\ell|)$, at most 100 iterations; if the decrement is still above round-off level a few extra steps are taken while they reduce the norm of the gradient. The covariance matrix is $(-H)^{-1}$ at the end. When Newton does not reach a maximum (no maximum exists, the Hessian is singular), \nwrem{} falls back to Mata's \texttt{optimize()}: the BFGS method from up to eight starting values (the first fixed, the others random), followed by Newton steps from the result, and, if BFGS fails, the Nelder--Mead method; the seed (\opt{seed()}) matters only there. If that fails as well, the command stops with return code 430.

\paragraph{Mata and the plugin.} Mata does the bookkeeping: it reads and sorts the events, writes the model into the frame, calls the plugin and runs the Newton loop. The plugin does the work that costs time. There is no slower fallback for the likelihood in Mata: a platform without a plugin cannot run \nwrem{} (return code 601).

\section{What the plugin computes}\label{sec:plugin}

The log rate of the pair $(j,k)$ at an event is a sum over the effects of $\theta_p x_p(j,k)$ (equation~(\ref{eq:rate})). Every statistic is the sum of a part of the sender, a part of the receiver and a sparse dyadic part that is nonzero only for the pairs that have been in contact (the inertia, recency and triadic effects), for the nonzero cells of a covariate network, or for the pairs of the previous event (participation shifts). The normalizer of the ordinal likelihood, $\sum_{(j,k)}\exp(\theta^{\top}x)$ over the risk set, is then the product of the sum over senders and the sum over receivers (without the diagonal), plus a correction over the sparse dyads only. The cost per event is therefore of the order $n$ plus the number of pairs in contact so far, and not $n^2$. The counts, the recency lists, the weighted two-path and shared-partner counts and the list of pairs in contact are updated after every event and do not depend on $\theta$. The sums that leave out the diagonal are formed from prefix and suffix sums, so that no cancellation arises.

When $\exp(\cdot)$ would overflow, a normalizer would underflow, or the sparse correction would cancel more than about six digits, the plugin evaluates that likelihood with a dense $O(n^2)$ evaluator instead, which is slower and equally exact. The fixed effects are always evaluated densely. The plugin carries no state between calls: the statistics do not depend on $\theta$, and everything that depends on the model is in the frame. The C source (\texttt{rem\_sim.c}, with StataCorp's plugin interface, which is not part of \nwrem{}'s own code) is part of the package (Section~\ref{sec:install}); the plugin does not depend on the version or the flavor of Stata (BE, SE, MP), only on the operating system and the processor. The plugin uses one processor core.

\paragraph{Precision.} The Hessian is exact, not a numerical derivative; the standard errors are those of the exact observed information. The estimates are the maximum of the likelihood to the round-off of the arithmetic, which is why they differ in the fifth digit or earlier from those that an optimizer with default tolerances reports.

\section{Speed}\label{sec:speed}

We timed \nwrem{}~1.4.0 against \texttt{rem.dyad()} of \relevent{}~\releventver{} (R~4.6.0), with the ordinal likelihood and maximum likelihood, on simulated event streams of 200 to 20,000 events among 20 to 100 actors and on five sets of effects: S1, two degree effects (\opt{nodsnd nidrec}); S2, four degree effects; S3, four degree effects and two inertia effects (\opt{frpsndsnd frrecsnd}); S4, two degree effects and two recency effects (\opt{rsndsnd rrecsnd}); S5, two degree effects and three covariate effects (\opt{covsnd covrec covint}). \nwrem{} always computes the covariance matrix from the exact Hessian, so the like-for-like setting of \relevent{} is \texttt{hessian=TRUE}; its default is \texttt{hessian=FALSE}, and both were timed.

\paragraph{Machine and method.} Intel Core i7-12700K (20 logical processors), 125\,GB memory, Ubuntu 22.04.4, StataNow~19.5 MP (the fits use one core), R~4.6.0; one server without other load. The times of \nwrem{}~\pkgversion{} were measured on 16 November 2026. The times of \relevent{}~\releventver{} were measured on 15 November 2026 on the same machine with an earlier version of \nwrem{} (1.2.0), and not repeated, so the two programs were not timed in turn. Each fit ran in its own Stata or R process after an untimed warm-up fit; the timer covers the estimation command (for R, the estimation call including its preprocessing), and loading the data is not timed. The table shows medians of up to 7 repeats with a different seed in each (5, 3 or 1 in the largest cells of \relevent{}). Both programs received the same data and the same specification with their default control settings. Each fit had a limit of 600 seconds. The data are simulated by generator scripts in the repository of the core package (folder \texttt{bench}). Timings depend on the machine, the data and the effects. The first call in a Stata session takes a short time longer, because the Mata library and the plugin are loaded; this is not in the times.

{\footnotesize\begin{longtable}{@{}llrrrrr@{}}
\toprule
Set & Events, actors & \nwrem{} (s) & \relevent{}, Hessian (s) & \relevent{}, no Hessian (s) & \multicolumn{2}{c}{Ratio with / without}\\
\midrule\endhead
S1 & 200, 20 & 0.006 & 0.133 & 0.107 & 22 & 18\\
 & 1\,000, 30 & 0.009 & 1.104 & 0.774 & 123 & 86\\
 & 5\,000, 20 & 0.016 & 3.240 & 2.478 & 202 & 155\\
 & 5\,000, 50 & 0.024 & 16.3 & 11.8 & 679 & 491\\
 & 5\,000, 100 & 0.041 & 58.5 & 40.7 & 1\,428 & 992\\
 & 20\,000, 100 & 0.213 & 396.1 & 320.3 & 1\,859 & 1\,504\\
\addlinespace
S2 & 200, 20 & 0.010 & 0.292 & 0.155 & 29 & 16\\
 & 1\,000, 30 & 0.013 & 3.213 & 1.636 & 247 & 126\\
 & 5\,000, 20 & 0.039 & 8.795 & 5.306 & 226 & 136\\
 & 5\,000, 50 & 0.048 & 49.5 & 29.3 & 1\,032 & 611\\
 & 5\,000, 100 & 0.068 & 151.4 & 73.4 & 2\,227 & 1\,079\\
 & 20\,000, 100 & 0.261 & -- & 576.4 & n.r. & 2\,208\\
\addlinespace
S3 & 200, 20 & 0.013 & 7.012 & 3.428 & 539 & 264\\
 & 1\,000, 30 & 0.051 & 106.3 & 49.2 & 2\,085 & 965\\
 & 5\,000, 20 & 0.176 & 312.2 & 174.6 & 1\,774 & 992\\
 & 5\,000, 50 & 0.667 & -- & -- & n.r. & n.r.\\
 & 5\,000, 100 & 1.334 & -- & -- & n.r. & n.r.\\
 & 20\,000, 100 & 9.534 & -- & -- & n.r. & n.r.\\
\addlinespace
S4 & 200, 20 & 0.014 & 0.402 & 0.249 & 29 & 18\\
 & 1\,000, 30 & 0.042 & 4.333 & 2.694 & 103 & 64\\
 & 5\,000, 20 & 0.137 & 12.8 & 8.103 & 93 & 59\\
 & 5\,000, 50 & 0.599 & 84.1 & 62.0 & 140 & 104\\
 & 5\,000, 100 & 0.792 & 203.8 & 125.4 & 257 & 158\\
 & 20\,000, 100 & 6.568 & -- & -- & n.r. & n.r.\\
\addlinespace
S5 & 200, 20 & 0.012 & 0.558 & 0.326 & 46 & 27\\
 & 1\,000, 30 & 0.015 & 5.587 & 3.187 & 372 & 212\\
 & 5\,000, 20 & 0.040 & 15.2 & 9.390 & 379 & 235\\
 & 5\,000, 50 & 0.053 & 81.6 & 48.5 & 1\,540 & 916\\
 & 5\,000, 100 & 0.083 & 295.4 & 169.2 & 3\,560 & 2\,038\\
 & 20\,000, 100 & 0.352 & -- & -- & n.r. & n.r.\\
\bottomrule
\end{longtable}
}
\noindent A dash means that \relevent{} did not finish within 600 seconds, or was not run because it had timed out in a smaller cell of the same effect set (\emph{n.r.} in the last two columns).

\paragraph{Reading the table.} \nwrem{} was faster than \relevent{} in every cell in which both programs finished: the time of \relevent{} divided by the time of \nwrem{} ranged from 22 to 3,560 with the Hessian and from 16 to 2,208 without. The smallest ratios are with the recency effects (S4: 29 to 257 with the Hessian). Fits with several hundred events take a few milliseconds. \relevent{} did not finish most of the fits with 20,000 events, while \nwrem{} finished all of them in 0.2 to 9.5 seconds; the largest fit, S3 with 20,000 events and 100 actors, took 9.5 seconds. The estimates agree: the largest relative difference of a coefficient between the programs was $4.2\cdot10^{-3}$, in a cell where the optimizer of \relevent{} stops early; in the other cells they agree to $10^{-3}$ or better.

\paragraph{Where the time goes.} The degree and actor covariate effects (S1, S2, S5) cost little because their statistics are sums over senders and receivers; the effects that depend on pairs in contact (S3, S4, the triadic effects, covariate networks with many nonzero cells) cost time in proportion to the number of such pairs, and the fixed effects are evaluated for every pair at every event. For a large network with many pairs in contact, the time grows with the number of pairs that have occurred.
